% !TEX encoding = UTF-8
\documentclass[UTF8]{ctexart}

\usepackage[utf8]{inputenc}
\usepackage{graphicx}
\usepackage{geometry}
\geometry{a4paper}
\geometry{left=2.5cm,right=2.5cm,top=2.8cm,bottom=1.3cm}

\usepackage{booktabs}
\usepackage{array}
\usepackage{paralist}
\usepackage{verbatim}
\usepackage{subfig}
\usepackage{amsmath}
\usepackage{mathtools}
\usepackage{listings}
\usepackage[table]{xcolor}
\usepackage{lastpage}
\usepackage{url}

% Using hyperref for improved ref character
\usepackage[colorlinks,linkcolor=black,anchorcolor=black,
citecolor=black,CJKbookmarks=True]{hyperref}

% For picture drawing
\usepackage[all]{xy}

% For code inserting. Set features.
\lstset{
alsolanguage=matlab,
tabsize=4,
keepspaces=true,
numbers=left,
numberstyle=\tiny,
keywordstyle=\color{blue!70} \bfseries,
commentstyle=\color{red!50!green!50!blue!50},
frame=shadowbox,
breaklines,
showspaces=false,
showstringspaces=false,
showtabs=false,
rulesepcolor=\color{red!20!green!20!blue!20},
extendedchars=false,
escapeinside=``
}

% Set the font of page header
\usepackage{fancyhdr}
\pagestyle{fancy}
\lhead{传感器特性实验报告}
\chead{}
\rhead{Page \thepage/\pageref{LastPage}}
\cfoot{}
\rfoot{}
\lfoot{}

\usepackage{sectsty}

\usepackage[nottoc]{tocbibind}
\usepackage[titles,subfigure]{tocloft}
\renewcommand{\cftsecfont}{\rmfamily\mdseries\upshape}
\renewcommand{\cftsecpagefont}{\rmfamily\mdseries\upshape}

% Set number of ref to be relevent to section number
\renewcommand{\theequation}{\arabic{section}.\arabic{equation}}
\renewcommand{\thefigure}{\arabic{section}-\arabic{figure}}
\renewcommand{\thetable}{\arabic{section}-\arabic{table}}
\makeatletter
\@addtoreset{equation}{section}
\@addtoreset{figure}{section}
\@addtoreset{table}{section}
\makeatother

% Set the font of the reference
\bibliographystyle{unsrt}

% Define user\rq{}s color
\usepackage{colortbl}
\definecolor{lightgray}{gray}{.9}
\definecolor{thickgray}{gray}{.6}

\usepackage{multirow}

% 首行缩进
\usepackage{indentfirst}

% Set section numbering
\CTEXsetup[number={}]{part}
\renewcommand{\thepart}{}
\usepackage{titlesec}
\titleformat{\part}[block]{\color{blue}\huge\bfseries\filcenter}{}{1em}{}

%\usepackage{ulem}
%\usepackage{indentfirst}
%\setlength\textwidth{300.0pt}
%

% 重定义字体命令
\newcommand{\song}{\CJKfamily{song}}    % 宋体   (Windows自带simsun.ttf)
\newcommand{\fs}{\CJKfamily{fs}}        % 仿宋体 (华天字库htfs.ttf)
\newcommand{\kai}{\CJKfamily{kai}}      % 楷体   (华天字库htkai.ttf)
\newcommand{\hei}{\CJKfamily{hei}}      % 黑体   (Windows自带simhei.ttf)
\newcommand{\li}{\CJKfamily{li}}        % 隶书   (Windows自带simli.ttf)
\newcommand{\you}{\CJKfamily{you}}      % 幼圆体 (Windows自带simyou.ttf)
%%%  以上六种字体均为标准 GBK 字体, 包括 GBK 繁体字和一些不常用字, 推荐!!!

\newcommand{\xs}{\CJKfamily{xs}}
\newcommand{\shu}{\CJKfamily{shu}}      % 舒体   (方正字库fzstk.ttf)
%  \newcommand{\yourcommand}[参数个数]{内容}   [参数个数]这个是可选的。
%  例如  \newcommand{\you}{\CJKfamily{you}}  用\you 来代替 \CJKfamily{you} ，少输入很多字哦。
%字号设置
\newcommand{\chuhao}{\fontsize{42pt}{\baselineskip}\selectfont}
\newcommand{\xiaochuhao}{\fontsize{36pt}{\baselineskip}\selectfont}
\newcommand{\yihao}{\fontsize{28pt}{\baselineskip}\selectfont}
\newcommand{\xiaoyihao}{\fontsize{24pt}{\baselineskip}\selectfont}
\newcommand{\erhao}{\fontsize{21pt}{\baselineskip}\selectfont}
\newcommand{\xiaoerhao}{\fontsize{18pt}{\baselineskip}\selectfont}
\newcommand{\sanhao}{\fontsize{15.75pt}{\baselineskip}\selectfont}
\newcommand{\sihao}{\fontsize{14pt}{\baselineskip}\selectfont}
\newcommand{\xiaosihao}{\fontsize{12pt}{\baselineskip}\selectfont}
\newcommand{\wuhao}{\fontsize{10.5pt}{\baselineskip}\selectfont}
\newcommand{\xiaowuhao}{\fontsize{9pt}{\baselineskip}\selectfont}
\newcommand{\liuhao}{\fontsize{7.875pt}{\baselineskip}\selectfont}
\newcommand{\qihao}{\fontsize{5.25pt}{\baselineskip}\selectfont}
% \baselineskip | distance from bottom of one line of a paragraph to bottom of the next line.  基本行距
%  只有使用\selectfont命令之后，\fontzize{}{}的设置才能生效
%  可以用数字表示{11pt}:单倍行距

\begin{document}
%%%%%%%%%%%%%%%%%%%%%%%%%%% Content %%%%%%%%%%%%%%%%%%%%%%%%%%%%%


\tableofcontents
\clearpage

\begin{abstract}

\end{abstract}
%%%%%%%%%%%%%%%%%%%%%%%%%% Main %%%%%%%%%%%%%%%%%%%%%%%%%%%%%%%%%%
[video to text introdution]
\section{Introduction}
The task of automatically describing videos containing
rich and open-domain activities poses an important challenges
for computer vision and machine learning research.
It also has a variety of practical applications.

\section{Related works}
\subsection{Sequence learning task}

\subsection{Image captioning}
Automatically generating a natural language description of an image, a problem known as image captioning, has recently received a lot of attention in Computer Vision. The problem is interesting not only because it has important practical applications, such as helping visually impaired people see, but also because it is regarded as a grand challenge for image understanding which is a core problem in Computer Vision.

\subsection{Video captioning}


\subsection{Attention}
Attention mechanism is a popular topic in deep learning these years. It is loosely based on human visual attention mechanism, as we tend to focus on a certain region with “high resolution” while perceiving the surrounding image in “low resolution”, and then adjust the focal point over time. Attention is successfully used in machine translation\cite{bahdanau2014neural}, image captioning\cite{xu2015show}, and video captioning\cite{venugopalan2014translating} tasks (these are the very first work using attention in their tasks).

\begin{figure}[htbp]
\centering
\includegraphics[width=11cm]{resources/logo.jpg}
\caption{No caption}
\label{fig:fullarm}
\end{figure}


\bibliography{reference}

\end{document}
